

Cette partie du projet correspond à ce que nous avons appelé notre \textbf{mode spécial}. Il s’agit d’une interface utilisateur en ligne de commande, développée en Python, qui permet d’accéder facilement aux analyses de performances. Ce menu, accessible via le script \texttt{menu\_performances.py}, joue le rôle de front-end : il propose à l’utilisateur différentes options pour consulter les statistiques issues de ses données GPS, comme la distance parcourue, la vitesse moyenne, l’altitude, les calories brûlées, ou encore les pauses détectées.

\vspace{1em}

Le menu repose sur le script \texttt{performances.py}, qui automatise l’analyse d’un fichier CSV contenant l’enregistrement du trajet. Ce traitement comprend plusieurs étapes: le nettoyage des données brutes, le calcul de différents indicateurs de performance, et la génération de graphiques. Cela permet d’obtenir une vue d’ensemble du parcours et de l’effort physique fourni.

\vspace{1em}

L’objectif de ce module est à la fois d’évaluer la qualité du suivi GPS (précision, cohérence des données) et de proposer des mesures fiables sur l’activité physique. L’ensemble du traitement est pensé pour être automatisé et reproductible, afin de gagner en rigueur et en lisibilité, tout en restant simple d’utilisation grâce à l’interface.

\vspace{2em}

\subsection{Données brutes et prétraitement}

Les données GPS sont enregistrées sous forme de tableau CSV. Chaque ligne contient un timestamp, des coordonnées (latitude et longitude) ainsi qu’une altitude.

\vspace{1em}

Avant l’analyse, un prétraitement est effectué pour fiabiliser les données. Les lignes incorrectes sont supprimées, les champs sont convertis en format numérique, et le timestamp est décomposé en date et heure. L’heure est ensuite exprimée en secondes depuis minuit, ce qui permet de calculer précisément les durées entre deux points.

\vspace{1em}

Ce traitement prépare les données pour les étapes suivantes, en assurant une base cohérente et exploitable.

\vspace{2em}

\subsection{Indicateurs de performance}

Une fois les données nettoyées, le script calcule automatiquement une série d’indicateurs. Il identifie les points de départ et d’arrivée, puis calcule la distance totale à l’aide de la formule de Haversine. La durée du trajet est obtenue à partir des heures converties, ce qui permet ensuite de déterminer la vitesse moyenne et la vitesse maximale.

\vspace{1em}

D’autres informations comme l’altitude minimale, maximale et moyenne sont également extraites. Le script estime aussi les calories brûlées à partir du poids, de la taille, et de la vitesse de l’utilisateur, en utilisant un facteur MET adapté à l’intensité de l’effort.

\vspace{1em}

Des pauses sont détectées en repérant les vitesses proches de zéro sur une durée significative. L’heure et la position GPS de chaque pause sont enregistrées pour une meilleure compréhension du déroulement du parcours. Tous ces calculs sont réalisés sans saisie manuelle.

\vspace{2em}

\subsection{Graphiques de visualisation des données}

Trois graphiques principaux sont générés automatiquement pour représenter le parcours :

\vspace{0.5em}

\begin{itemize}
	\item L’altitude en fonction du temps, pour visualiser les phases de montée et de descente.
	\item La vitesse en fonction du temps, pour repérer les variations d’allure et les arrêts.
	\item L’altitude en fonction de la distance parcourue, pour une lecture plus concrète du relief.
\end{itemize}

\vspace{0.5em}

Ces représentations visuelles facilitent l’interprétation des données et permettent de repérer d’éventuelles anomalies, comme des valeurs extrêmes ou incohérentes.

\vspace{2em}

\subsection{Interprétation des indicateurs et qualité du suivi}

Les indicateurs calculés permettent d’évaluer la fiabilité du système. Ils montrent si les données GPS sont cohérentes avec le terrain, si les vitesses enregistrées sont réalistes, et si les pauses ont bien été identifiées. En parallèle, l’utilisateur obtient une estimation précise de l’effort physique fourni, ce qui peut être utile pour le suivi d’une activité ou d’un entraînement.

\vspace{1em}

L’ensemble du système a été conçu pour fournir une analyse complète, avec des résultats clairs, et accessibles grâce à un menu interactif simple à utiliser. Cela rend le traitement non seulement automatisé, mais aussi directement exploitable par toute personne souhaitant consulter ses performances.

\vspace{2em}

Ce module regroupe à la fois le traitement statistique des données GPS et une interface utilisateur permettant de naviguer facilement dans les résultats. Le menu en ligne de commande rend l’ensemble du système plus accessible et interactif, tout en conservant une logique de traitement rigoureuse et structurée. Cela représente un vrai point fort du projet, en associant analyse technique et utilisation pratique.

